%%%%%%%%%%%%%%%%%%%%%%%%%%%%%%%%%%%%%%%%%%%%%%%%%%%%%%%%%%%%
%% Lecture 08
%%%%%%%%%%%%%%%%%%%%%%%%%%%%%%%%%%%%%%%%%%%%%%%%%%%%%%%%%%%%

\lecture
{08}
{Wednesday, 2 September 2026}
{Introduction to Probability}
{Dr. Nishant Chandgotia}

%%%%%%%%%%%%%%%%%%%%%%%%%%%%%%%%%%%%%%%%%%%%%%%%%%%%%%%%%%%%
%% Independence and Expectations
%%%%%%%%%%%%%%%%%%%%%%%%%%%%%%%%%%%%%%%%%%%%%%%%%%%%%%%%%%%%

Suppose $X,Y\in L^1(\mathbb{P})$ and $X$ and $Y$ are independent.
Then
\(
    \mathbb{E}(XY)
    =
    \mathbb{E}(X)\mathbb{E}(Y).
\)

\begin{question}{Why does this make sense?}{Quesito-3}
    Why does \(XY \in L^{1}\left( \mathbb{P} \right)\)?
\end{question}

\begin{remark}
The basic idea is that if $X\sim\mu$ and $Y\sim\nu$ are
independent, then

\[
    \mathbb{P}(X\in A,Y\in B)
    =
    \mathbb{P}(X\in A)\mathbb{P}(Y\in B)
    =
    \mu(A)\nu(B).
\]

Hence,
\(
    (X,Y)\sim\mu\times\nu.
\)

Therefore,

\[
\begin{aligned}
    \mathbb{E}(|XY|)
    &=
    \mathbb{E}(|X||Y|) \\
    &=
    \int_{\mathbb{R}^2}
    |x||y|\,d(\mu\times\nu)(x,y).
\end{aligned}
\]

By Fubini's theorem,

\[
\begin{aligned}
    \mathbb{E}(|XY|)
    &=
    \left(
        \int_{\mathbb{R}} |x|\,d\mu(x)
    \right)
    \left(
        \int_{\mathbb{R}} |y|\,d\nu(y)
    \right) \\
    &=
    \mathbb{E}(|X|)\mathbb{E}(|Y|)
    <\infty.
\end{aligned}
\]

The same argument applied to $XY$ gives

\[
    \mathbb{E}(XY)
    =
    \mathbb{E}(X)\mathbb{E}(Y).
\]

This also works when $X,Y\geq 0$.

Similarly, if
\(
    X_1,X_2,\ldots,X_n\in L^1(\mathbb{P})
\)
are independent, then

\[
    \mathbb{E}
    \left(
        \prod_{i=1}^n X_i
    \right)
    =
    \prod_{i=1}^n \mathbb{E}(X_i).
\]
\end{remark}

%%%%%%%%%%%%%%%%%%%%%%%%%%%%%%%%%%%%%%%%%%%%%%%%%%%%%%%%%%%%
%% Covariance
%%%%%%%%%%%%%%%%%%%%%%%%%%%%%%%%%%%%%%%%%%%%%%%%%%%%%%%%%%%%

\begin{definition}{Covariance}{covariance}
For random variables $X,Y\in L^2(\mathbb{P})$, the covariance of
$X$ and $Y$ is defined by

\[
    \operatorname{Cov}(X,Y)
    =
    \mathbb{E}
    \left[
        (X-\mathbb{E}(X))(Y-\mathbb{E}(Y))
    \right].
\]

Equivalently,

\[
    \operatorname{Cov}(X,Y)
    =
    \mathbb{E}(XY)
    -
    \mathbb{E}(X)\mathbb{E}(Y).
\]
\end{definition}

\begin{remark}
If $X$ and $Y$ are independent, then
\(
    \operatorname{Cov}(X,Y)=0.
\)

The converse is not true in general.
\end{remark}

%%%%%%%%%%%%%%%%%%%%%%%%%%%%%%%%%%%%%%%%%%%%%%%%%%%%%%%%%%%%
%% Counterexample
%%%%%%%%%%%%%%%%%%%%%%%%%%%%%%%%%%%%%%%%%%%%%%%%%%%%%%%%%%%%

\begin{example}{Uncorrelated but Dependent Random Variables}{uncor-but-dep-rv}
Consider random variables $X,Y$ taking values in
$\{-1,0,1\}$ with joint probability table

\[
\begin{array}{c|ccc}
    & Y=1 & Y=0 & Y=-1 \\ \hline
X=1  & 0 & a & 0 \\
X=0  & b & c & b \\
X=-1 & 0 & a & 0
\end{array}
\]

where
\(
    2a+2b+c=1,
    \qquad
    a,b,c\geq 0.
\)

Clearly this is a probability distribution.

We have
\(
    \mathbb{P}(X=1,Y=1)=0,
\)
but
\(
    \mathbb{P}(X=1)\mathbb{P}(Y=1)>0
\)
whenever $a,b>0$. Hence $X$ and $Y$ are not independent.

On the other hand,
\(
    \mathbb{E}(X)=0,
    \qquad
    \mathbb{E}(Y)=0,
\)

and therefore

\[
    \operatorname{Cov}(X,Y)
    =
    \mathbb{E}(XY)
    -
    \mathbb{E}(X)\mathbb{E}(Y)
    =
    0.
\]

Thus zero covariance does not imply independence.
\end{example}

%%%%%%%%%%%%%%%%%%%%%%%%%%%%%%%%%%%%%%%%%%%%%%%%%%%%%%%%%%%%
%% Questions
%%%%%%%%%%%%%%%%%%%%%%%%%%%%%%%%%%%%%%%%%%%%%%%%%%%%%%%%%%%%

\begin{question}{Distribution of a Transformation}{distribution-transformation}
If
\(
    X\sim\operatorname{Unif}[-1,1],
\)
what is the distribution of
\(
    X^3?
\)
\end{question}

\begin{question}{Bernoulli Random Variables}{bernoulli-random-variables}
Suppose
\(
    X,Y\sim\operatorname{Ber}.
\)

What conditions on $X$ and $Y$ imply independence?
\end{question}

\begin{question}{Jointly Normal Random Variables}{jointly-normal}
Suppose $X$ and $Y$ are jointly normal random variables.

Does
\(
    \operatorname{Cov}(X,Y)=0
\)
imply that $X$ and $Y$ are independent?
\end{question}

\begin{question}{Bernoulli Variables and Covariance}{bernoulli-covariance}
Suppose
\(
    X\sim\operatorname{Ber}(p),
    \qquad
    Y\sim\operatorname{Ber}(q).
\)

If
\(
    \operatorname{Cov}(X,Y)=0,
\)
does it follow that $X$ and $Y$ are independent?
\end{question}

%%%%%%%%%%%%%%%%%%%%%%%%%%%%%%%%%%%%%%%%%%%%%%%%%%%%%%%%%%%%
%% Inequalities
%%%%%%%%%%%%%%%%%%%%%%%%%%%%%%%%%%%%%%%%%%%%%%%%%%%%%%%%%%%%

\section{Inequalities}

\begin{definition}{Markov's Inequality}{markov-inequality}
Let $X\geq 0$ be a random variable and let $a>0$. Then

\[
    \mathbb{P}(X\geq a)
    \leq
    \frac{\mathbb{E}(X)}{a}.
\]
\end{definition}

\begin{remark}
The main idea behind Markov's inequality is the following.
Let
\(
    \varphi:\mathbb{R}\to\mathbb{R}
\)
be a measurable function such that
\(
    \varphi\geq 0.
\)

For $A\in\mathcal{F}$,

\[
    \mathbb{E}
    \left[
        \varphi(X)\mathbf{1}_{\{X\in A\}}
    \right]
    \leq
    \mathbb{E}[\varphi(X)].
\]

This simple inequality has many important consequences.
\end{remark}

%%%%%%%%%%%%%%%%%%%%%%%%%%%%%%%%%%%%%%%%%%%%%%%%%%%%%%%%%%%%
%% Chebyshev
%%%%%%%%%%%%%%%%%%%%%%%%%%%%%%%%%%%%%%%%%%%%%%%%%%%%%%%%%%%%

\begin{definition}{Chebyshev's Inequality}{chebyshev-inequality}
Let
\(
    \mu=\mathbb{E}(X),
    \qquad
    \sigma^2=\operatorname{Var}(X),
    \qquad
    \lambda>0.
\)

Then

\[
    \mathbb{P}
    \left(
        |X-\mu|\geq\lambda\sigma
    \right)
    \leq
    \frac{1}{\lambda^2}.
\]
\end{definition}

\begin{remark}
Let, 
\(
    Y=(X-\mu)^2.
\)
Then
\(
    \mathbb{E}(Y)
    =
    \operatorname{Var}(X)
    =
    \sigma^2.
\)

Therefore, by Markov's inequality,

\[
\begin{aligned}
    \mathbb{P}
    \left(
        |X-\mu|\geq\lambda\sigma
    \right)
    &=
    \mathbb{P}
    \left(
        Y\geq\lambda^2\sigma^2
    \right) \\
    &\leq
    \frac{\sigma^2}
    {\lambda^2\sigma^2} \\
    &=
    \frac{1}{\lambda^2}.
\end{aligned}
\]
\end{remark}

%%%%%%%%%%%%%%%%%%%%%%%%%%%%%%%%%%%%%%%%%%%%%%%%%%%%%%%%%%%%
%% Chernoff
%%%%%%%%%%%%%%%%%%%%%%%%%%%%%%%%%%%%%%%%%%%%%%%%%%%%%%%%%%%%

\begin{definition}{Chernoff Bound}{chernoff-bound}
For $t>0$, define the moment generating function

\[
    M(t)=\mathbb{E}(e^{tX}).
\]

Then

\[
    \mathbb{P}(X\geq a)
    \leq
    \frac{M(t)}{e^{ta}}.
\]
\end{definition}

\begin{remark}
The Chernoff bound follows immediately by applying Markov's
inequality to the nonnegative random variable $e^{tX}$:

\[
\begin{aligned}
    \mathbb{P}(X\geq a)
    &=
    \mathbb{P}(e^{tX}\geq e^{ta}) \\
    &\leq
    \frac{\mathbb{E}(e^{tX})}{e^{ta}} \\
    &=
    \frac{M(t)}{e^{ta}}.
\end{aligned}
\]
\end{remark}

%%%%%%%%%%%%%%%%%%%%%%%%%%%%%%%%%%%%%%%%%%%%%%%%%%%%%%%%%%%%
%% Coupon Collector Problem
%%%%%%%%%%%%%%%%%%%%%%%%%%%%%%%%%%%%%%%%%%%%%%%%%%%%%%%%%%%%

\section{Coupon Collector Problem}

\begin{question}{Coupon Collector Problem}{coupon-collector-problem}
Suppose there are $n$ different types of cards. Each time a card is
picked, each of the $n$ types is obtained independently and uniformly.

How many cards do we need to draw to ensure that we have obtained all
$n$ types of cards?
\end{question}

Consider the probability space
\(
    \Omega
    =
    \left(
        \{1,2,\ldots,n\}^{\mathbb{N}},
        \mathcal{B},
        \operatorname{Unif}^{\otimes\mathbb{N}}
    \right).
\)

Let
\(
    X_n(\omega)
    =
    \#\{\text{cards picked to obtain the full collection}\}.
\)

Equivalently,

\[
    X_n(\omega)
    =
    \min
    \left\{
        m:
        \{\omega_1,\omega_2,\ldots,\omega_m\}
        =
        \{1,\ldots,n\}
    \right\}.
\]

\begin{theorem}{Coupon Collector Bound}{coupon-collector-bound}
For every $\varepsilon>0$,

\[
    \mathbb{P}
    \left(
        1-\varepsilon
        <
        \frac{X_n}{n\log n}
        <
        1+\varepsilon
    \right)
    \longrightarrow 1
    \qquad\text{as }n\to\infty.
\]
\end{theorem}

\begin{proof}
We prove the upper bound. It is enough to show that

\[
    \mathbb{P}
    \left(
        \frac{X_n}{n\log n}>1+\varepsilon
    \right)
    \longrightarrow 0
    \qquad\text{as }n\to\infty.
\]

Equivalently,

\[
    \mathbb{P}
    \left(
        X_n>(1+\varepsilon)n\log n
    \right)
    \longrightarrow 0.
\]

Let $Y_k$ denote the number of card types missing after $k$
draws. Thus,

\[
    Y_k(\omega)
    =
    n-
    \left|
        \{\omega_1,\omega_2,\ldots,\omega_k\}
    \right|.
\]

We want to show that if
\(
    k>(1+\varepsilon)n\log n,
\)
then $Y_k$ is very likely to be zero.

For each $i\in\{1,\ldots,n\}$ and each $k\geq 1$, define

\[
    S_{i,k}
    =
    \begin{cases}
        1,
        & \text{if card type $i$ is not picked by time $k$,}\\
        0,
        & \text{otherwise.}
    \end{cases}
\]

Then
\(
    Y_k
    =
    \sum_{i=1}^n S_{i,k}.
\)

For a fixed $i$,
\(
    \mathbb{E}(S_{i,k})
    =
    \left(1-\frac{1}{n}\right)^k,
\)
since the probability that type $i$ is not selected in one draw is
$1-\frac1n$.

Therefore,

\[
    \mathbb{E}(Y_k)
    =
    n
    \left(1-\frac{1}{n}\right)^k.
\]

Now take ~ ~ ~
\(
    k=(1+\varepsilon)n\log n.
\)

By Markov's inequality,

\[
\begin{aligned}
    \mathbb{P}
    \left(
        X_n>(1+\varepsilon)n\log n
    \right)
    &=
    \mathbb{P}
    \left(
        Y_{(1+\varepsilon)n\log n}\geq 1
    \right) \\
    &\leq
    \mathbb{E}
    \left[
        Y_{(1+\varepsilon)n\log n}
    \right] \\
    &=
    n
    \left(
        1-\frac{1}{n}
    \right)^{(1+\varepsilon)n\log n}.
\end{aligned}
\]

Using
\(
    1-\frac1n\leq e^{-1/n},
\)
we obtain

\[
\begin{aligned}
    n
    \left(
        1-\frac1n
    \right)^{(1+\varepsilon)n\log n}
    &\leq
    n
    \exp
    \left(
        -\frac{(1+\varepsilon)n\log n}{n}
    \right) \\
    &=
    n e^{-(1+\varepsilon)\log n} \\
    &=
    n^{-\varepsilon}.
\end{aligned}
\]

Hence,

\[
    \mathbb{P}
    \left(
        X_n>(1+\varepsilon)n\log n
    \right)
    \leq
    n^{-\varepsilon}
    \longrightarrow 0.
\]

Thus,

\[
    \mathbb{P}
    \left(
        \frac{X_n}{n\log n}>1+\varepsilon
    \right)
    \longrightarrow 0.
\]

This proves the desired upper bound.
\end{proof}

%%%%%%%%%%%%%%%%%%%%%%%%%%%%%%%%%%%%%%%%%%%%%%%%%%%%%%%%%%%%
%% End of Lecture
%%%%%%%%%%%%%%%%%%%%%%%%%%%%%%%%%%%%%%%%%%%%%%%%%%%%%%%%%%%%

\lectureend



\section{Quiz, 4 September 2026}

\begin{quiz}

\begin{enumerate}

    \item It is equivalent to a question discussed in class.
    What is it?

    \item If
    \(
        X\sim\operatorname{Geo}(p),
    \)
    prove that
    \(
        \mathbb{E}(X)=\frac{1}{p},
        \qquad
        \operatorname{Var}(X)\leq\frac{1}{p^2}.
    \)

    \item Let $X_k^{(n)}$ be the time taken for $k$ bins to be
    filled. Prove that
    \(
        S_k^{(n)}
        =
        X_{k+1}^{(n)}-X_k^{(n)}
        \sim
        \operatorname{Geo}(p_k),
    \)
    where
    \(
        p_k=1-\frac{k}{n},
    \)
    and that the random variables
    \(
        S_0^{(n)},S_1^{(n)},\ldots,S_{n-1}^{(n)}
    \)
    are independent.

    \item Prove that
    \(
        \mathbb{E}\left(X_n^{(n)}\right)
        =
        n\sum_{m=1}^{n}\frac{1}{m}.
    \)

    \item Prove that
    \(
        \operatorname{Var}\left(X_n^{(n)}\right)
        \leq cn^2
    \)
    for some constant $c$.

    \item Prove that
    \(
        \frac{\mathbb{E}\left(X_n^{(n)}\right)}
        {n\log n}
        \longrightarrow 1
        \qquad\text{as }n\to\infty.
    \)

    \item Use Chebyshev's inequality to conclude that
    \(
        \frac{
            X_n^{(n)}
            -
            n\displaystyle\sum_{m=1}^{n}\frac{1}{m}
        }{
            nf(n)
        }
        \longrightarrow 0
        \qquad\text{in probability as }f(n)\to\infty.
    \)
\end{enumerate}

\end{quiz}
